\subsection{The sigmoid as a smoothed-notch approximation}
\label{app:notch}

The calibrated model of Section~\ref{sec:model} treats the registration
threshold as a true notch and characterises its analytic objects---the
dominated region and the marginal buncher---exactly. For the forward solution at
the level of two million firms, I operationalise the notch with a smooth
(logistic) approximation to the tax schedule. This appendix sets out that
approximation and establishes when it is faithful: the smooth schedule
reproduces the notch---including its dominated region---when its steepness $k$ is
calibrated to the Kleven--Waseem width, and degenerates towards a smoothed
\emph{kink} only when $k$ is mis-set. I calibrate $k$ accordingly and report all
exact analytic objects from the true notch, so the smoothing affects only the
numerics of the forward solution, not the paper's quantitative claims.

\subsubsection*{The smoothing device}

The exact notch makes the effective average rate a step function of turnover,
non-differentiable at $T^{*}$. To obtain a smooth schedule, I replace the step
by a logistic effective-rate schedule,
\begin{equation}
  \tau(y_i) \;=\; \tau_{\max}\,\frac{1}{1+e^{-k(y_i - T^{*})}},
  \qquad
  \frac{\partial \tau(y_i)}{\partial y_i}
  \;=\; k\,\tau(y_i)\!\left(1-\frac{\tau(y_i)}{\tau_{\max}}\right),
  \label{eq:tau-smooth}
\end{equation}
where $k>0$ governs the steepness: as $k\to\infty$ the sigmoid converges
pointwise to the exact step at $T^{*}$, and its derivative peaks at $T^{*}$ and
decays away from it. The corresponding optimality condition equates marginal
after-tax revenue $(1-\tau(y_i))-y_i\,\partial\tau/\partial y_i$ to marginal cost.
I do not solve this condition as a root-find: with $\alpha\approx0.989$ the
near-constant-returns specification makes a continuous interior optimum
ill-conditioned, so the operational solver maximises after-tax profit directly
over a bounded grid (described below), which does not rely on the optimality
condition having a well-defined interior root.

\subsubsection*{Fidelity of the smooth approximation}

Equation~\eqref{eq:tau-smooth} smooths the \emph{average} rate $\tau(y_i)$
applied to the firm's entire turnover, so after-tax revenue is
$R(y_i)=(1-\tau(y_i))y_i$. Whether this reproduces the institutional notch
depends on the steepness $k$. Under the true notch, after-tax revenue drops
discretely by $\tau_{\max}T^{*}=\pounds17{,}000$ at the threshold; under the
smooth schedule the drop is spread over a transition whose width falls as $k$
rises, and whether $R(y_i)$ decreases over a neighbourhood of $T^{*}$---the
structural analogue of the output-dominated region---holds for steep $k$ and
fails for shallow $k$, where the schedule collapses to a smoothed kink that
merely perturbs the marginal rate locally.

The consequences for the dominated region are stark.
Table~\ref{tab:notch-vs-sigmoid} reports the width of the locally
output-dominated region implied by the sigmoid (the band where marginal
after-tax revenue $(1-\tau)-y\tau'$ turns negative) against the exact
Kleven--Waseem width $a=T^{*}\tau_{\max}/(1-\tau_{\max})=\pounds21{,}250$
of~\eqref{eq:dominated}. Only at a steepness close to $k=0.001$ does the implied
band approach the exact width; at the larger $k=0.01$ used in earlier code the
sigmoid captures barely a tenth of it, and for very small $k$ the schedule is
monotone and exhibits \emph{no} dominated region at all---a pure kink. The full
$k$ grid and its reconciliation with the empirical excess span are reported in
Appendix~\ref{app:inference}.

\begin{table}[t]
  \centering
  \caption{Sigmoid steepness $k$ and the implied output-dominated region,
    compared to the exact Kleven--Waseem notch width
    $a=T^{*}\tau_{\max}/(1-\tau_{\max})=\pounds21{,}250$ at
    $T^{*}=\pounds85{,}000$, $\tau_{\max}=0.20$. The true notch is
    $k$-independent.}
  \label{tab:notch-vs-sigmoid}
  \begin{tabular}{lcc}
    \hline
    Schedule & Dominated-region width & Character \\
    \hline
    Sigmoid, $k=0.001$      & $\approx\pounds19.7$k     & notch-like \\
    Sigmoid, $k=0.01$       & $\approx\pounds2.5$k      & narrow, near-kink \\
    Sigmoid, $k=0.0001$     & none (monotone $R$)       & pure kink \\
    True notch (this paper) & $\pounds21{,}250$ (exact) & notch \\
    \hline
  \end{tabular}
\end{table}

Two points reinforce the message. First, the sigmoid's dominated-region width is
almost entirely a function of $k$, which was not calibrated to any data moment in
earlier drafts and took inconsistent values across paper and code ($k=0.01$
versus $k=0.001$); any bunching prediction read off the sigmoid inherits this
arbitrariness. Second, for the questions this paper asks---the size of the
bunching band, the location of the marginal buncher, and the revenue
consequences of moving the notch---these discrepancies are first-order, which is
why the main text models the notch exactly.

\subsubsection*{Recommended use and reconciliation}

I use the smooth schedule of~\eqref{eq:tau-smooth} as the effective-rate
schedule in the forward firm-by-firm re-optimisation, but the firm's optimum is
located by \emph{bounded grid maximisation of after-tax profit}, not by an
interior root-find. For each firm I evaluate after-tax profit on a fine turnover
grid over a bounded neighbourhood ($\pm\pounds20$k on a \pounds200 step) of its
observed turnover and select the profit-maximising point. I adopt direct profit
maximisation precisely because the near-constant-returns specification
($\alpha\approx0.989$) makes the interior optimum degenerate, so the global
maximiser is well-defined while an interior root is not. The
$\pm\pounds20$k window is itself a modelling assumption, motivated by the
documented localisation of the realised response within roughly \pounds20{,}000
of the threshold (Section~\ref{sec:results_structural}); I note that this choice
mechanically bounds how far firms can move, so the ``localised, partial'' response
partly reflects the window rather than being a free output of the exercise.

I adopt two further safeguards throughout. (i) I calibrate $k$ to the exact
dominated-region width $a=T^{*}\tau_{\max}/(1-\tau_{\max})$ ($k\approx0.001$,
which reproduces the band to within about $7\%$;
Table~\ref{tab:notch-vs-sigmoid}) rather than asserting it, and use a single value
consistently. (ii) I report every exact analytic object---the
$\sim\pounds21$k dominated region and the marginal buncher of
Section~\ref{ssec:buncher}---from the true notch of Section~\ref{sec:model}, not
from the smoothing. Consistent with this posture, I headline the taper's robust
revenue effect and its analytic efficiency mechanism (the removal of the dominated
region and the discrete wedge); the \emph{simulated} change in the near-threshold
turnover distribution depends on both $k$ and the re-optimisation window and is
not reported as a magnitude (Section~\ref{sec:results_taper}).
